\documentclass[12pt,a4paper]{article}
% Times New Roman 12pt, single spacing: matches the SAS CU Hackathon Round-2 report format.
\usepackage[T1]{fontenc}
\usepackage{mathptmx}
\usepackage[margin=1.8cm,top=1.6cm,bottom=1.6cm]{geometry}
\usepackage{booktabs,tabularx,array}
\usepackage{xcolor}
\usepackage[framemethod=default]{mdframed}
\usepackage{enumitem}
\usepackage{hyperref}

\definecolor{sasblue}{HTML}{0766D1}
\newmdenv[linecolor=sasblue,linewidth=1.2pt,backgroundcolor=sasblue!6,
  innertopmargin=6pt,innerbottommargin=6pt,innerleftmargin=8pt,innerrightmargin=8pt]{objbox}
\newcolumntype{L}[1]{>{\raggedright\arraybackslash}p{#1}}
\newcolumntype{Y}{>{\raggedright\arraybackslash}X}
\setlist{nosep,leftmargin=*}
\pagestyle{empty}

\begin{document}

\begin{center}
  {\Large\bfseries Analytics Objective}\\[2pt]
  {\large D2S Bharat: From Data-Science Job Demand to Skills, Success and Training Decisions}\\[4pt]
  {\small Team Code5urge \quad$\cdot$\quad SAS CU Hackathon / Build for Bharat \quad$\cdot$\quad
  \textit{Draft for team agreement}}
\end{center}

\noindent\textbf{Problem context.}
Institutions that prepare data professionals must decide what to teach with limited budget,
trainers and seats. They currently lack evidence on three questions: which skills the market
\emph{demands and pays for}, which skills are associated with \emph{early-career growth}, and which
personal traits are associated with \emph{success among senior} data scientists. The SAS-provided data covers all three:
job postings, company-level openings and salaries, and two outcome-labelled employee datasets.

\vspace{6pt}
\begin{objbox}
\noindent\textbf{Analytics objective.}
To identify the technical skills and personality traits most strongly associated with
labour-market demand, pay and on-the-job success for data-science professionals, and to
translate this evidence into a budget-constrained, explainable training plan that tells an
institution \emph{what to teach first}.
\end{objbox}

\vspace{4pt}
\noindent\textbf{Research questions.}\\[2pt]
{\footnotesize
\begin{tabularx}{\linewidth}{@{}L{1.9cm}Y L{3.4cm} L{3.6cm}@{}}
\toprule
\textbf{RQ} & \textbf{Question} & \textbf{Data} & \textbf{Methods} \\
\midrule
RQ1 Demand \& pay & Which skills and roles are most demanded, and which are associated with
higher salary bands and experience requirements? &
Analytics Jobs (14,840 postings after cleaning); DataScience Jobs (1,602 rows, 93,005 openings) &
Cleaning, skill normalisation to ESCO/O*NET, descriptive statistics, non-parametric tests \\
\addlinespace
RQ2 Junior growth & Which of the five measured skill areas are associated with a high vs.\ low
salary hike among junior data scientists? & JDS Skill Traits (139) &
Group tests, correlation, regularised logistic regression, shallow decision tree,
repeated cross-validation \\
\addlinespace
RQ3 Senior success & Which Big Five traits are associated with high vs.\ low success among
senior, customer-facing data scientists? & SDS Personality Traits (161) & As RQ2 \\
\addlinespace
RQ4 Decision & Given RQ1--RQ3, which training interventions should an institution prioritise
under budget and capacity limits? & Evidence from RQ1--RQ3 &
Constrained optimisation (OR-Tools CP-SAT), sensitivity analysis \\
\bottomrule
\end{tabularx}}

\vspace{6pt}
\noindent\textbf{Scope and assumptions.}
Only the SAS-provided datasets are analysed; ESCO and O*NET serve solely as a skill vocabulary.
Findings are associations, not causal effects. JDS and SDS are small samples, so model
complexity is limited deliberately: with five predictors each, the smaller outcome class gives
13.2 (JDS) and 15.2 (SDS) events per predictor, a favourable ratio, while model complexity
stays deliberately constrained. We use
L2-regularised logistic regression and decision trees of depth $\le 3$, repeated stratified
k-fold cross-validation, and bootstrap confidence intervals for effects. Salaries are analysed
as band midpoints.

\vspace{4pt}
\noindent\textbf{Stakeholders and expected outputs.}
\begin{itemize}
  \item \textbf{Institutions and curriculum designers:} a ranked, evidence-backed list of skills
        to teach, and an optimised training plan with its cost--impact trade-offs.
  \item \textbf{Students:} which skills and traits are linked to pay, growth and success.
  \item \textbf{Employers and policy/CSR programmes:} where supply and demand diverge, by role and city.
\end{itemize}

\vspace{4pt}
\noindent\textbf{Success criteria.}
Each RQ is answered with a quantified finding; predictive models (RQ2, RQ3) are judged against
a majority-class baseline using cross-validated accuracy and AUC; and the RQ4 plan states which
constraints bind and what extra budget would change.

\end{document}
